% !TeX root = main.tex
\section{Document Conventions}

\subsection{Update Project}

The phrase \emph{update project} means:
\begin{enumerate}
\item Read \texttt{main.tex} and every active file that it includes.
\item Compare every model document included by \texttt{main.tex} against its
current \texttt{<model>Dynamics} implementation.
\item Analyze which document sections are out of date.
\item Present the proposed document changes for review before modifying any
document file.
\item After approval, update the affected document sections.
\item Add a new entry to \texttt{ProblemEvolution.tex} when the approved work
changes the dynamics model or records a completed investigation.
\item Do not modify \texttt{<model>DynamicsFlowTikz.tex} unless that specific
change is separately reviewed and approved.
\item Present a proposed git commit message.
\end{enumerate}

\subsection{Problem Evolution}
\label{sec:problem-evolution}

The intent of this section is to follow the path of dynamics evolution.
Understanding that path makes the design issues and nomenclature visible.  It
also helps future development because the current model contains assumptions
about rules such as equal-radius particles, wall particles, and source-owned
contact state.  When those rules change, it should be clear where they came
from and why they existed.

It shall take the form:
\begin{verbatim}

\subsection{<Problem Statement>}

\subsubsection{<previous form>}
\subsubsection{<current form>}

\subsubsection{difference outcome}
	one of :full solution, particle solution, no solution

\subsubsection{<Analysis>}

\end{verbatim}

\subsection{git message}

\begin{verbatim}
The git message will be of the format :
Current status:
	- <what we changed since last commit>
	- <etc>
Code Changes:
	- <short description of code changes this session>
	- <etc>
Next: <What we will do next>
\end{verbatim}

\subsection{Document Authority}

The compiled document rooted at \texttt{main.tex} is the single dynamics
document.  Its included files divide the document into maintainable sections;
they are not separate competing dynamics specifications.

The dynamics implementations are <model>Dynamics.  Documentation
must describe that implementation and its current architecture.  Historical
architecture names and retired implementations do not belong in the active
document.

\subsection{Model Documents}

Each maintained dynamics model has a corresponding
\texttt{<model>DynamicsSource.tex} file that documents that model's source
flow, equations, model-specific state, and implementation architecture.  Every
maintained model source document must be included by \texttt{main.tex}.

Shared document conventions, dynamics conventions, programming conventions,
and problem evolution remain in their corresponding shared document files.
Retired model source documents must be removed from \texttt{main.tex} or moved
under the old-document directory.

\subsection{Retiring Dynamics}

Retiring \texttt{<model>Dynamics} means:
\begin{enumerate}
\item Remove the model import and registry entry from
\texttt{SimulationRunner\_HSVPanel.py}.
\item Move its model-specific Python files under
\texttt{python/base/old/<model>Dynamics}.
\item Remove its source document from \texttt{main.tex} and move any retained
model-specific documents under \texttt{docs/Dynamics/old}.
\item Ensure \texttt{ParticleUtil.cfg} does not select the retired model.
\end{enumerate}

\subsection{Generated Source Flowcharts}

For dynamics work, the sequential stage chain should be
documented with a generated flowchart when the source structure becomes hard to
hold in memory.  The flowchart is not a substitute for narrow functions; it is a
debugging and navigation aid that makes the stage order visible.

The current source-flow generator is:

\begin{verbatim}
python/base/<model>DynamicsFlowChart.py
\end{verbatim}

It emits DOT, SVG, PDF, and PNG artifacts under
\texttt{docs/Dynamics/generated}.  The DOT file preserves the graph source, the
SVG can carry interactive links, and the PDF is used as the stable LaTeX figure.
The source-flow chart links mostly to code.  The system-flow chart links mostly
to documentation because it explains loops, decision trees, storage boundaries,
and GLSL-shaped algorithm stages.  The generators should be updated when
function names, stage order, decision points, or major source handoffs change.

This preserves the same programming rule used by the code: the primary flow is
explicit, generated documentation follows that flow, and helpers are shown as
supporting calls rather than hidden control-flow owners.
